\section{Consequences of the further proof attempts}
\label{next-status}

The new results change the research program. One proposed strengthening
is false in the model as defined, a proposed vacuum reference is not a
minimizer, and the screened lattice intermediate theorem can be proved
without a spectral gap. These statements did not settle the original
unscreened many-body comparison. That historical open problem is now
resolved in the positive-effective-mass regime by the primary paper
\emph{Vacuum-relative stability of cutoff QED}.

\subsection{Three outcomes and their exact scope}

\begin{enumerate}
\item \textbf{Point sources prevent the proposed ultraviolet uniformity.}
At fixed \(m,\alpha,L\) satisfying Theorem~\ref{uv-theorem},
\[
 -C\alpha\Lambda_f
 \le E^{\rm HF}_{\Lambda,L}(0,1)-E^{\rm HF,\vac}_{\Lambda,L}
 \le m-c_\alpha\Lambda_f .
\]
The lower bound is Theorem~\ref{thm:mf}; the upper bound is the explicit
neutral Slater construction. Thus the linear cutoff order is necessary
for this one-nucleus problem. The failure occurs with \(Z=1\), one
neutralizing electron, and at most one coherent pair. It is not
supercritical one-particle Dirac collapse or an effect requiring
infinitely many real particles.

In particular, control of the anomalous charge density by a kinetic
estimate with a cutoff-independent point-source error cannot complete
the previously proposed argument. The packet construction exploits a
term linear in the pair rotation, while kinetic and self-consistent
Coulomb costs begin quadratically. Keeping a finite cutoff is consistent
with the proved theorem. Removing it requires a genuinely different
formulation or a renormalization analysis. Independently calibrated
per-species source masses are legitimate state- and position-independent
parameters; this bare-reference counterexample does not refute that
separately normalized question.

\item \textbf{Static leading-order matching does not identify the vacuum.}
At \(\delta\mu^2=2e^2\chi(0)\),
\[
 E^{\vac,(\delta\mu)}_{\Lambda,L}
 \le E^{\rm prod,\vac}_{\Lambda,L}
 \le E_\gamma(\delta\mu)-cL^3 .
\]
The proof retains exchange and uses exact pointwise zero charge density.
It is therefore stronger than the failure of a convenient absorption
inequality. The actual product-state vacuum energy must be subtracted
in any subsequent relative argument.

Corollary~\ref{pv-exchange-threshold} gives a necessary condition on a
bare mass for the free sea to remain a minimizer. It is not a construction
of a massless physical photon, a proof of a particular ordered phase, or
a sufficient stability criterion. The extensive lower bound for a
positive bare mass remains valid.

\item \textbf{Screened lattice relative stability holds for correlated states.}
For the explicit lattice model of Section~\ref{sl-yukawa},
\[
 E_\nu-E_0^{\vac}\ge-C_\mu Z_{\max}K,
 \qquad
 |c_\nu-c_0|\le2C_\mu Z_{\max}K .
\]
Here the reference is the actual interacting minimum in the neutral
sector, and the source state is in its different, exactly neutralizing
charge sector. The proof handles this difference explicitly. It permits
arbitrary pair coherence and correlated many-body states; it does not
replace the many-body infimum by a Hartree--Fock infimum.

The theorem uses an absolutely summable source potential. Its constant
contains \(4\pi/\mu^2\), and therefore does not survive removal of
screening by the same argument. The bound on the correlation-energy
difference is a consequence of the two-sided energy comparison, not
a small correlation-energy expansion.
\end{enumerate}

\subsection{What remains before a full-QED theorem}

\paragraph{A vacuum comparison at fixed cutoff with Coulomb tails.}
The primary target at the time of this historical section was
\[
 E_{\Lambda,L}(R,Z)-E_{\Lambda,L}^{\vac}\ge-C_\Lambda K
\]
for the unscreened, pair-producing many-body model in its stated
large-box range. The primary paper now proves this in its specified
positive-effective-mass regime, including quantized photons.
The raw \(L^5\) vacuum construction does not disprove
this subtraction. The new ultraviolet counterexample is a statement
about the mean-field reference and changing cutoff; it also does not
disprove the fixed-cutoff many-body comparison.
For a chosen model with extensive vacuum energy, a useful next theorem
would control the energy change caused by external sources relative
to that vacuum with an error proportional to their number. The
screened lattice result supplies a test case but no estimate uniform
in the screening length.

\paragraph{A suitable reference for the photon problem.}
The exchange calculation requires replacing the free sea as a presumed
minimizer at the static tuning. At each fixed volume and cutoff,
the product-state variational problem has a minimizer: after photon
minimization it is a continuous function on the compact Grassmannian
of exact-charge projections. This elementary existence statement gives
neither a thermodynamic vacuum nor its coercivity. A thermodynamic
construction and localized response estimates are separate tasks.
If a gauge-covariant lattice model is adopted instead, gluing must
retain the Gauss-law boundary-flux sectors described earlier.

\paragraph{Source and ultraviolet normalization.}
A finite nuclear radius, a different reference sea, and independently
specified source-renormalization conditions are different possible
models, not consequences of the present proof. The point-source
counterexample establishes why the original fixed-reference
cutoff-independent claim cannot simply be recovered by sharpening its
constants. Any proposed modification must state its energy convention
before the stability inequality is tested. No physical mass or charge
matching has been supplied by the variational calculations here.

\paragraph{Low-energy reduction and equilibrium matter.}
The new relative lattice estimate is an energy bound; it does not
provide a pair-sector spectral gap. It therefore supplies neither of
the missing configuration-uniform Feshbach inputs automatically.
One must still establish a suitable pair-sector resolvent bound and
an extensive estimate on virtual-pair corrections before transferring
nonrelativistic stability back to the relativistic Hamiltonian.
Thermodynamic limits, atomic binding, and crystalline equilibrium also
remain distinct questions. The three results proved here remove
specific incorrect routes and establish a nonperturbative intermediate
theorem; they do not manufacture the missing full-QED argument.
